% Source: https://artofproblemsolving.com/community/c3h2530554_what_exactly_is_conditional_probability
% Stable identity: aops-topic-2530554-q01
\textbf{(a)} If there are $N$ equally likely outcomes, then
\[P(A\mid B)=\dfrac{P(A\cap B)}{P(B)}=\dfrac{\dfrac{|A\cap B|}{N}}{\dfrac{|B|}{N}}=\dfrac{|A\cap B|}{|B|}.\]
Knowing that $B$ happened shrinks the set of possible outcomes to $B$, whose outcomes are still equally likely; $P(A\mid B)$ is the probability of $A$ in this smaller sample space.

\textbf{(b)} There are $36$ equally likely outcomes. No $6$ appears in $5\cdot5=25$ of them, so $P(A)=\dfrac{11}{36}$. The outcomes with sum $8$ are $(2,6),(3,5),(4,4),(5,3),(6,2)$, and two of them, $(2,6)$ and $(6,2)$, contain a $6$. So
\[P(A\mid B)=\dfrac25,\qquad P(B\mid A)=\dfrac{|A\cap B|}{|A|}=\dfrac{2}{11}.\]
Since $P(A\mid B)=\dfrac25\ne\dfrac{11}{36}=P(A)$, the events are \textbf{not} independent: knowing the sum is $8$ makes a $6$ more likely.

\[\boxed{P(A)=\dfrac{11}{36},\quad P(A\mid B)=\dfrac25,\quad P(B\mid A)=\dfrac{2}{11},\quad\text{not independent}}\]
